\documentclass{article}

\usepackage[utf8]{inputenc}
\usepackage[turkish, shorthands = :!]{babel}
\usepackage[T1]{fontenc}
\usepackage{lmodern}

\usepackage{fancyhdr, lipsum, mathtools}

\usepackage{amsthm} % Typesetting theorems (AMS style)

\theoremstyle{plain} % varsayılan
\newtheorem{prop}{Önerme}
\newtheorem{lem}{Önsav}
\newtheorem{conj}{Sanı}
\newtheorem{cor}{Sonuç}
\newtheorem{thm}{Teorem}
\newtheorem*{thm*}{Teorem} % numaralandırılmamış teorem
\newtheorem*{thm1*}{Filancanın Teoremi} % özel adlandırılmış teorem

\theoremstyle{definition}
\newtheorem{defn}{Tanım}
\newtheorem{ex}{Örnek}

\theoremstyle{remark}
\newtheorem*{rem}{Not}

\newtheoremstyle{solution}{}{}{}{}{\itshape}{.}{ }{} % ``ispat'' ortamının ``çözüm'' sürümü gibi
\theoremstyle{solution}
\newtheorem*{soln}{Çözüm}

%\newtheorem{conj}{Conjecture}
%\newtheorem{cor}{Corollary}
%\newtheorem{lem}{Lemma}
%\newtheorem{prop}{Proposition}
%\newtheorem{thm}{Theorem}

% numaralar biçimden bağımsız tek bir sıra izlesin istersek:
%\newtheorem{prop}{Önerme}
%\newtheorem{lem}[prop]{Önsav}
%\newtheorem{conj}[prop]{Sanı}
%\newtheorem{cor}[prop]{Sonuç}
%\newtheorem{thm}[prop]{Teorem}

\addto\captionsturkish{
	\renewcommand{\proofname}{İspat} % "kanıt" yerine "ispat"
}

\renewcommand{\headrulewidth}{0 pt}
\pagestyle{fancy}
\fancyhf{}
\fancyhead[C]{\thepage}

\begin{document}

\lipsum[1][1-6]

\begin{defn}
Bu bir tanımdır.
\end{defn}

\lipsum[1][7-12]

\begin{ex}
Bu bir örnektir.
\end{ex}

\begin{soln}
Bu bir çözümdür.
\end{soln}

\begin{ex}
Bu başka bir örnektir.
\end{ex}

\lipsum[1][7-12]

\begin{rem}
Bu bir nottur.
\end{rem}

\lipsum[1][13-18]

\begin{lem} \label{lem:1}
Bu bir önsavdır.
\end{lem}

Önsav~\ref{lem:1}'den dolayı...

\begin{prop}
Bu bir önermedir.
\end{prop}

\lipsum[2][1-6]

\begin{thm}
Bu bir teoremdir.
\end{thm}

\begin{proof}
Bu bir ispattır.
\end{proof}

\lipsum[2][7-12]

\begin{thm*}
Numaralandırılmamış bir teorem.
\end{thm*}

\begin{thm1*}
Özel adlandırılmış bir teorem.
\end{thm1*}

\end{document}